% Part: second-order-logic
% Chapter: metatheory
% Section: compactness

\documentclass[../../../include/open-logic-section]{subfiles}

\begin{document}

\olfileid{sol}{met}{com}

\olsection{د دويمې درجې منطق کمپکت نۀ دے}

\begin{explain}
د جملو سټ~$\Gamma$ هغه وخت \emph{په متناهي ډول د صدق وړ} بلئ چې هر
متناهي فرعي سټ ئې د صدق وړ وي۔ د لومړۍ درجې منطق دا خاصيت لري چې
که د !!{sentence}s سټ~$\Gamma$ په متناهي ډول د صدق وړ وي، نو پخپله
هم د صدق وړ دے۔ دا خاصيت \emph{کمپکتوالی} بلل کېږي۔ يوه معادله بڼه
ئې د منطقي لازم والي په هکله ده: که $\Gamma \Entails !A$، نو د کوم
متناهي فرعي سټ~$\Gamma_0 \subseteq \Gamma$ دپاره لا له وړاندې
$\Gamma_0 \Entails !A$۔ په دې بڼه دا د بشپړتيا د قضيې سمدستي پايله
ده: که $\Gamma \Entails !A$، نو د بشپړتيا له مخې $\Gamma \Proves !A$۔
خو يو !!{derivation} د~$\Gamma$ له متناهي شمېر !!{sentence}s څخه
ګټه اخيستلے شي۔

کمپکتوالی د دويمې درجې منطق دپاره نۀ صدق کوي۔ د دويمې درجې د
!!{sentence}s داسې سټونه شته چې په متناهي ډول د صدق وړ دي، خو
پخپله د صدق وړ نۀ دي؛ او داسې سټونه هم شته چې له هغو څخه~$!A$
منطقي لازم راځي، خو له هېڅ متناهي فرعي سټ څخه ئې~$!A$ نۀ راځي۔
\end{explain}


\begin{thm}
\ollabel{thm:sol-not-compact}
د دويمې درجې منطق کمپکت نۀ دے۔
\end{thm}
\textbf{د ژباړې د نسخې سمون (OLSOL-004):} سرچينه دلته د مخکينۍ
ناپرېکړه‌وړتيا د قضيې عين ليبل بيا کاروي؛ د کمپکتوالي قضيې ته
ځانګړے ليبل ورکړل شو۔

\begin{proof}
په ياد ولرئ چې
\[
\fn{Inf} \ident \lexists[u][(\lforall[x][\lforall[y][(\eq[u(x)][u(y)]
      \lif \eq[x][y])]] \land \lexists[y][\lforall[x][\eq/[y][u(x)]]])]
\]
په !!a{structure} کښې هله او يوازې هله صدق کوي چې تعريف ساحه ئې
نامتناهي وي۔ دې $!A^{\ge n}$ هغه جمله وي چې وايي تعريف ساحه لږ تر
لږه $n$ !!{element}s لري، لکه:
\[
!A^{\ge n} \ident \lexists[x_1][\dots\lexists[x_n][
    \bigwedge_{1 \le i < j \le n}\eq/[x_i][x_j]]].
\]
\textbf{د ژباړې د نسخې سمون (OLSOL-005):} د سرچينې په بېلګه کښې
يوازې ځينې جوړې بېلې ښودل شوې وې؛ لږ تر لږه د ټاکلي شمېر بېلابېل
غړي غوښتل د ټولو بېلو جوړو شرط غواړي۔
د !!{sentence}s لاندې سټ وګورئ:
\[
\Gamma = \{\lnot \fn{Inf}, !A^{\ge 1}, !A^{\ge 2}, !A^{\ge 3}, \dots\}.
\]
دا په متناهي ډول د صدق وړ دے، ځکه چې د هر متناهي فرعي سټ~$\Gamma_0
\subseteq \Gamma$ دپاره داسې $k \ge 1$ غوره کولے شو چې له~$k$ څخه د لوړو
اندازو هېڅ جمله په دغه فرعي سټ کښې نۀ وي۔
\textbf{د ژباړې د نسخې سمون (OLSOL-006):} په سرچينه کښې دا
متناهي‌والي شرط په تېروتنې سره ټول نامتناهي سټ ته منسوب شوے و؛
دلته شرط يوازې پر متناهي فرعي سټ لګېږي۔ که $\Domain{M}$
$k$ !!{element}s ولري، نو $\Sat{M}{\Gamma_0}$۔ خو $\Gamma$ د صدق وړ
نۀ دے: که $\Sat{M}{\lnot \fn{Inf}}$، نو $\Domain{M}$ بايد متناهي وي،
مثلاً د اندازې~$k$۔ بيا $\Sat/{M}{!A^{\ge k+1}}$۔
\end{proof}

\begin{prob}
داسې سټ~$\Gamma$ او !!a{sentence}~$!A$ مثال ورکړئ چې
$\Gamma \Entails !A$ وي، خو د هر متناهي فرعي سټ~$\Gamma_0 \subseteq
\Gamma$ دپاره $\Gamma_0 \Entails/ !A$ وي۔
\end{prob}


\end{document}
